\section{Conclusion}
\label{sec:conclusion}
HMRS makes the state used to manage documents available to the agents that
retrieve them. A catalogue and shelves expose membership and descriptions;
document and segment views expose text with source references. These views
make three checks possible in the tested workflows: is the current document
searchable, did the agent reach the intended collection, and what text
reached the evidence pack?

The experiments show why these checks matter. A lifecycle probe located
stale entries in the index. The highest target-string count in
the frozen search configurations came from A0h, which combined document
text and prioritised chunk windows; reaching the correct collection alone
was not sufficient. The fixed-pool replay shows why admission scores should
be tested separately from ranking. The next test is whether these practices
improve claim-level support on independent questions when agents have the
same search and reading budgets.
